% cap04/4.1_patron_headless.tex
\section{Segregaci\'on y Desacoplamiento (Modelo Headless)}

La arquitectura de \textit{Democra} repudia por completo el dise\~no tradicional de aplicaciones monol\'iticas, donde el c\'odigo visual y la l\'ogica transaccional coexisten y se ejecutan en el mismo servidor. En su lugar, documenta la adopci\'on de un paradigma de \textbf{Arquitectura en M\'ultiples Capas con Desacoplamiento Extremo (\textit{Headless})}. Este enfoque segrega de forma absoluta las responsabilidades del sistema, aislando quir\'urgicamente el \textit{Frontend} del \textit{Backend}.

\subsection{El \textit{Frontend} en el \textit{Edge} (\textit{Edge Computing})}

El \textit{Frontend}, construido \'integramente en \textit{React.js} bajo el patr\'on \textit{Single Page Application (SPA)} y la filosof\'ia \textit{Mobile First}, no reside en ning\'un servidor tradicional. Es transpilado y empaquetado como micro-assets est\'aticos (archivos \textit{.js, .css, .html}). Esta arquitectura dicta que estos \textit{assets} sean despachados a trav\'es de \textit{Content Delivery Networks (CDN)} de alta disponibilidad en el \textbf{Edge (\textit{Edge Computing})}.

Al posicionar el c\'odigo visual en nodos \textit{Edge} (como \textit{AWS CloudFront} o \textit{Vercel}), se reduce dr\'asticamente la latencia de descarga. El navegador del usuario (sea un dispositivo m\'ovil andino o una estaci\'on de trabajo) descarga la SPA desde el servidor m\'as cercano. Una vez descargada, la interfaz (React) no reside en el servidor, sino que instila su estado enteramente dentro de la memoria temporal (\textit{Virtual DOM}) del navegador.

\subsection{El \textit{Backend} como Orquestador Sem\'antico Puro}

Por su parte, el \textit{Backend} de \textit{Democra}, soportado por \textit{Node.js} y \textit{Express}, se despoja de cualquier prop\'osito de renderizado. Su \textbf{\'unico prop\'osito inquebrantable} es fungir como un orquestador algor\'itmico. Est\'a dise\~nado exclusivamente para:
\begin{enumerate}
    \item \textbf{Recibir verbos sem\'anticos HTTP:} (\textit{GET, POST, PUT, DELETE}) desde la SPA.
    \item \textbf{Contrastar heur\'isticamente los tokens:} (\textit{JSON Web Tokens, JWT}) para autenticaci\'on y autorizaci\'on.
    \item \textbf{Sanitizar la entrop\'ia de las cargas:} (\textit{Payloads JSON}) inyectadas por los clientes, neutralizando cualquier vector de ataque.
    \item \textbf{Traducir al protocolo de base de datos:} orquestando la comunicaci\'on con \textit{PostgreSQL}.
\end{enumerate}

Este \textit{Backend} \textbf{no genera un s\'olo p\'ixel} de interfaz gr\'afica. Responde exclusivamente con flujos de datos puros (\textit{JSON}).

\subsection{Tuber\'ias Cifradas (TLS / HTTPS)}

Para garantizar la integridad de la informaci\'on que viaja a trav\'es del \textit{Edge} hacia el orquestador central, la segregaci\'on se sella con \textbf{tuber\'ias cifradas}. El tr\'afico \textit{Frontend-Backend} viaja encriptado bajo el protocolo \textit{Transport Layer Security (TLS)} encapsulado en \textit{HTTPS}, asegurando la inviolabilidad de las credenciales y tokens en tr\'ansito.
